\section{Discussion}
\label{sec:discussion}

\subsection{What the search is actually buying}
\label{sec:disc-mechanism}

The ablation of Section~\ref{sec:eval-ablation} is the most useful negative result in this
work: budget-matched uniform sampling of the schema grammar matched or exceeded the
evolutionary search on the held-out half of every cell. Read carelessly, that says the
optimizer does not work. Read correctly, it locates the contribution precisely. The gain
comes from two things, and neither of them is the search heuristic. The first is the
\emph{grammar}: a space that contains single, nested, and conditional designs at all
parameterizations, so that the best available answer is inside it. The second is
\emph{measured selection on the real query stream}, which is what distinguishes a good
answer from a plausible one. Any procedure that samples the grammar densely enough and
measures honestly will find comparable schemas.

This matters for adoption more than a better optimizer would. A practitioner does not have
to implement an evolutionary loop, tune its mutation rates, or trust its convergence. They
have to enumerate a design space and measure candidates against a captured trace. The
evolutionary machinery in our implementation is retained because it is cheap and
occasionally converges sooner, not because the results depend on it.

The corollary is a limit on how far the present results generalize. We can claim that
composing indexes under a measured objective beats a fixed default. We cannot claim that
our particular traversal of the design space is efficient, and we do not.

\subsection{Why the winners look the way they do}
\label{sec:disc-locality}

The schemas the search selects share a bias toward shallow, wide structures, and the reason
is an implementation property rather than a geometric one. Regressing measured latency on
the structural counters separates the two cost terms in a traversal: a visited node costs
roughly \pendingnum{$50$--$130$~ns} in our CPU implementation, against roughly
\pendingnum{$1$--$3$~ns} for a tested point. The per-point figure is what a float distance
test should cost. The per-node figure is far above a bounds check, and it is dominated by
memory: nodes are individually allocated, a large kd-tree spans tens of megabytes, and
visiting a few dozen scattered nodes in it is a few dozen cache misses.

An index whose nodes were contiguous and cache-resident would shift that ratio, and with
it the optimum. Designs that touch more nodes to test fewer points would become cheaper
relative to designs that do the opposite, and the search would move toward them. Our
conclusions about \emph{which} schema wins are therefore conditional on a traversal whose
node visits are expensive. The conclusion that \emph{measurement should decide} is not: a
different implementation changes the answer, which is precisely the argument for searching
per instance rather than legislating a default.

We treat this openly as the largest optimization opportunity left in the system, and as a
caution against reading our selected schemas as statements about spatial data structures in
general. They are statements about spatial data structures as implemented here, measured
under the workload given.

\subsection{Where schema search does not pay}
\label{sec:disc-failure}

Three regimes emerged in which the honest output of the search is not a nested design.

\paragraph{Small clouds.} Below roughly ten million points a tuned single primitive wins.
The index fits in cache, the traversal is short, and the extra indirection of a second level
costs more than the pruning it buys. The search returns the single primitive, which is the
behavior we asked of it in Section~\ref{sec:introduction}, but it means the method has no
value at that scale beyond confirming the default.

\paragraph{A hyper-engineered incumbent.} The triangle-mesh boundary study
(Section~\ref{sec:eval-boundary}) is the sharpest case. Against a SAH-built compressed wide
BVH on a real scanned model, every wrapper and every nested schema loses, and the optimizer
converges to the incumbent. Decades of engineering effort concentrated on one structure for
one query type is not something a generic composition search should be expected to beat.
The synthetic axis-aligned scene, where a searched grid nearly doubles throughput, shows the
other side of the same coin: the search wins where no such incumbent exists. Mixed,
dense point workloads are that setting; primary-ray traversal of triangle meshes is not.

\paragraph{Parity cells.} In cells where several designs are within measurement noise, the
selected winner varies between repeated searches. We report winner families rather than
exact schemas there. A practitioner in such a cell gains nothing from tuning, and the
useful output of the search is the negative finding itself: any of these designs will do.

\subsection{GPU support is a coverage gap, not a speed result}
\label{sec:disc-gpu}

The CUDA path exists to expose backend-specific limitations rather than to claim a device
speedup, and the honest summary is that its coverage is partial. Not every primitive in the
grammar has a CUDA builder, so a nested schema can be partially supported; the framework
records support status per row and refuses to rank across backends, because a schema that
silently fell back to a CPU level would otherwise look like a fast GPU schema. Device memory
is the harder limit. Positions alone consume 12~bytes per point, so a
billion-point cloud fills most of a 16~GB device before any nodes exist, and the GPU arm of
our evaluation is capped well below the top of the scale ladder.

Nothing here says GPU indexing is a dead end. It says that a fair CPU-versus-GPU comparison
is a different paper, needing device-resident query streams and a memory budget the present
protocol does not model, and that mixing the two backends in one ranking would produce a
number that means nothing.

\subsection{Per-instance tuning versus a general selector}
\label{sec:disc-transfer}

The economics of this method depend on amortization. A search costs far more than a single
query, and is worth running only because the workloads we target issue millions to billions
of them against a stable distribution. Two results bound how far one search can be spread.

Transfer across clouds of the same morphology is \emph{safe but not optimal}: the
cross-cloud experiment of Section~\ref{sec:eval-transfer} places a transferred schema
within a few percent of that cloud's own best, while a wrongly fixed default costs
multiples. Reusing a same-morphology schema is therefore a reasonable default and re-tuning
recovers the remainder. Transfer across scale behaves similarly: rankings obtained on a
small subsample of a cloud largely survive at full scale, which is what makes searching at
low fidelity and confirming at high fidelity affordable.

Both results point toward a learned selector that predicts a schema from cloud and workload
features rather than searching per instance, and the framework already exports its
measurements in the form such a selector would need. We stop short of that claim here. The
evidence supports interpolation within a morphology and within a scale ladder; it does not
support a general predictor, and the honest statement of the boundary is that per-instance
measurement remains the fallback whenever the target is outside the range where transfer
has been demonstrated.

\subsection{Limitations}
\label{sec:disc-limitations}

Beyond the threats to validity in Section~\ref{sec:eval-threats}, four limitations bear on
how these results should be used.

\emph{Static indexes.} The framework builds an index over a fixed point set. Insertion,
deletion, and streaming updates are outside the schema space, so nothing here speaks to
workloads where the cloud changes under the index. Update cost is one of the classical axes
of spatial-index design, and its absence narrows the design space we search.

\emph{Coordinate precision.} Positions are stored as single-precision offsets from a
per-cloud origin. Over a scene tens of kilometres across this leaves a floor of a few
millimetres, which is immaterial for surveys at centimetre density but becomes a real
fraction of the point spacing for a merged scene that spans kilometres while holding
centimetre-resolution detail. Where a scene approaches that regime we report the floor
alongside the measurement rather than assuming exactness, and exactness checks against
double-precision baselines are interpreted against it.

\emph{Trace provenance.} The pipeline traces are derived from the geometry of the stages
rather than instrumented inside a third-party library. For the self-query stages this is
exact by construction --- the stage queries every point, and the trace is every point --- but
the neighborhood radii are set from a measured mean point spacing, so the trace encodes our
estimate of the scale at which each stage operates rather than a specific library's
parameter choice.

\emph{Dataset licensing.} Two of the mobile-mapping scenes are distributed under
non-commercial licences, one of them additionally prohibiting derivative distribution. We
report measurements and figures for them, and the derived multi-scale tiers for that scene
cannot be redistributed with the rest of the benchmark. Reproduction there requires
rebuilding the tiers from the original download.

\section{Conclusions}
\label{sec:conclusions}

We have treated point-cloud index selection as a physical-design problem and shown that it
can be answered by measurement rather than by convention. Given a cloud and a captured query
trace, a search over a grammar of single, nested, and conditional index schemas selects a
replayable design that outperforms both the best tuned single primitive and hand-designed
nested baselines on held-out queries, on scans of every morphology we tested, once the cloud
is large enough for traversal cost to matter.

Three findings seem to us more durable than any individual speedup. First, the cost of
choosing wrongly is larger than the cost of choosing well is small: a fixed default index
carried onto the wrong cloud costs multiples of the best available choice, which is the
practical argument for measuring at all. Second, the benefit comes from the design space
combined with honest measurement, not from the optimizer that traverses it --- uniform
sampling does as well as evolutionary search, which makes the method far easier to adopt
than to invent. Third, a measured search is trustworthy exactly because it is willing to
return a negative result: against a hyper-engineered incumbent it converges to that
incumbent, and on small clouds it returns a single primitive.

The workloads that motivate this are the ones where the index is an inner loop rather than a
convenience: per-point geometry-processing stages, and the neighborhood queries inside
hierarchical point networks. In both, the query distribution is known before the work
starts and is issued a very large number of times, so a one-off search is amortized to
nothing. That regime is common, growing, and poorly served by a fixed default.

The clearest next step follows from Section~\ref{sec:disc-locality}. Node visits in our
traversal are dominated by cache misses, and the selected schemas reflect that. Making node
storage contiguous would not merely make the system faster; it would change which designs
win, and re-running the same search over the same clouds would then measure how much of the
present design space is an artifact of memory layout rather than of geometry. That
experiment is the natural sequel, and the framework is built to answer it by measurement
rather than by argument.
